\documentclass[letterpaper,12pt]{article}
%\usepackage{harvard}
%\usepackage{subfigure}
%\usepackage[active]{srcltx}
\usepackage{graphicx}
\usepackage{pdflscape}
\usepackage{amsmath}
\usepackage{amsthm}
\usepackage{lscape}
\usepackage{enumerate}
\usepackage{float}
\usepackage{grffile}
\usepackage{relsize}
\usepackage{tablefootnote}
\usepackage{footnote}
\usepackage{sectsty}
\usepackage{xcolor}
\usepackage{caption}
\usepackage{color}
\usepackage{natbib}
\usepackage{multirow}
\usepackage{bbm}
\usepackage[T1]{fontenc}
\usepackage[multiple]{footmisc}
\usepackage{appendix}
\usepackage{times}
\usepackage{relsize}
\usepackage{caption}
\usepackage{newtxtext,newtxmath}
\usepackage{enumitem}
\usepackage{tikz}


%%%%%%%%%%%%%%%%%%%%%%%%%%%%%%%%%%%%%%%%%%%%%%%%%%%%%%%%%%%%
%%%   Title Page Information
%%%%%%%%%%%%%%%%%%%%%%%%%%%%%%%%%%%%%%%%%%%%%%%%%%%%%%%%%%%%
\title{\vspace{-1in}ECON311, Winter 2026: Final Exam}
\author{Leandro Freylejer}
\date{}

%%%%%%%%%%%%%%%%%%%%%%%%%%%%%%%%%%%%%%%%%%%%%%%%%%%%%%%%%%%%
%%%   Double and Single Space Commands
%%%%%%%%%%%%%%%%%%%%%%%%%%%%%%%%%%%%%%%%%%%%%%%%%%%%%%%%%%%%
\newlength{\defbaselineskip}
\setlength{\defbaselineskip}{\baselineskip}
\newcommand{\setlinespacing}[1]{\setlength{\baselineskip}{#1 \defbaselineskip}}
\newcommand{\doublespacing}{\setlength{\baselineskip}{1.3 \defbaselineskip}}
\newcommand{\singlespacing}{\setlength{\baselineskip}{1\defbaselineskip}}
\renewcommand{\thesubsection}{\arabic{subsection}}
\newtheoremstyle{dotlessS}{}{}{\itshape}{}{\color{dg}\bfseries}{}{ }{}
\theoremstyle{dotlessS}
\newtheorem{theorem}{Theorem}
\newtheorem{lemma}{Lemma}
\newtheorem{proposition}{Proposition}
\newtheorem{assumption}[theorem]{Assumption}
\newtheorem{corollary}{Corollary}
\newtheorem{definition}{Definition}
\newenvironment{example}[1][Example]{\begin{trivlist}
\item[\hskip \labelsep {\bfseries #1}]}{\end{trivlist}}
\newenvironment{remark}[1][Remark]{\begin{trivlist}
\item[\hskip \labelsep {\bfseries #1}]}{\end{trivlist}}

\newcommand\Pitem{%
  \addtocounter{enumi}{-1}%
  \renewcommand\theenumi{\arabic{\enumi}'}%
  \item%
  \renewcommand\theenumi{\arabic{enumi}}%
}

%%%%%%%%%%%%%%%%%%%%%%%%%%%%%%%%%%%%%%%%%%%%%%%%%%%%%%%%%%%%
%%%   Set Margins Here
%%%%%%%%%%%%%%%%%%%%%%%%%%%%%%%%%%%%%%%%%%%%%%%%%%%%%%%%%
\textwidth=6.5in
\textheight=9in
\oddsidemargin=0.10in
\evensidemargin=0.10in
\topmargin=-0.5in
%\parskip=3pt

%%%%%%%%%%%%%%%%%%%%%%%%%%%%%%%%%%%%%%%%%%%%%%%%%%%%%%%%%%%%%
%%%   Actual Document Begins Here
%%%%%%%%%%%%%%%%%%%%%%%%%%%%%%%%%%%%%%%%%%%%%%%%%%%%%%%%%%%%%
\begin{document}
\pagenumbering{arabic}
\maketitle
\noindent There are two parts in this exam. Part I has five questions and Part II has two questions for a total of 83 points. You must show all your work to get full marks. You are allowed a calculator. Good Luck!!!
\section{Part I: Short-Answer Questions}

\subsection*{Question I: International Transmission of Monetary Policy (6 points)}
Suppose that the Federal Reserve in the US decides to implement contractionary monetary policy due to high inflation ($\uparrow i^{US}$). Assuming that the Canadian economy is at the potential output level, use the IS-MP model and the Phillips Curve to analyse how this policy affects the Canadian economy in the short run. To answer this question, you do not need to derive all the equations of the model. However, you need to carefully discuss the effect of the policy on the value of the Canadian Dollar relative to the US Dollar and show the effect of this shock graphically (in both, the IS-MP and Phillips Curve graphs). Suppose that the Bank of Canada wants to return the economy to its potential output level, how would it use domestic monetary policy to counter this shock?     

\subsection*{Question II: The Golden-Rule Level of Capital (6 points)}
Using the Solow Model studied in class, briefly explain what the golden-rule level of capital is. In your answer, make sure to explain what the golden-rule level of capital maximizes and use the Solow diagram to illustrate your answer.

\subsection*{Question III: Romer Model}
Consider the following Romer economy: 
\begin{eqnarray*}
\begin{array}{ll}
\text{Production Function:} & Y_t=A_t  \: L_{Y_t} \\[1.4ex]
\text{Productivity:} & A_{t+1}=A_t+\overline{z} \: A_t \: L_{A_{t}}  \\[1.4ex]
\text{Allocation of labour:} & L_{A_{t}}=\overline{I} \: \overline{L}, \: \: \:  \: \: L_{Y_{t}}=(1-\overline{I}) \: \overline{L}  \\[1.4ex]
\end{array}
\end{eqnarray*}
\begin{enumerate}[label=(\alph*)]
\item (4 points) Derive the rate of growth of output per person in the long-run equilibrium ($g_y$)
\item (5 points) Use the previous equation to explain the trade-off between long-run growth and the short-run level of consumption in this model. In your answer make sure to include a graph with output per person on the vertical axis (using a ratio scale) and time on the horizontal axis. Label your graph carefully, including slopes.
\end{enumerate} 




\subsection*{Question IV: Solow Model (6 points)}
Refer to Figure 1 when answering the following question.
\begin{figure}[H]
\centering
% Three graphs in one row: K, C, Y -- all increasing to new higher steady state (decrease in depreciation rate)
\begin{minipage}[t]{0.31\textwidth}
\centering
\begin{tikzpicture}[font=\small, scale=0.75]
  \draw[->] (0,0) -- (6.2,0) node[right] {$t$};
  \draw[->] (0,0) -- (0,4.8) node[above] {$K_{\text{\itshape (ratio scale)}}$};
  \draw[thin, dashed] (2,0) node[below] {$t_{shock}$} -- (2,4.5);
  \draw[thick] (0.2,1.5) -- (2,1.5);
  \draw[thick] (2,1.5) .. controls (2.8,2.5) and (5.0,3.2) .. (5.9,3.2);
  \draw[thin, dash dot] (2,3.2) -- (6.2,3.2);
  \node[left] at (0,1.5) {$K^{*}$};
  \node[left] at (0,3.2) {$K^{**}$};
\end{tikzpicture}
\end{minipage}
\hfill
\begin{minipage}[t]{0.31\textwidth}
\centering
\begin{tikzpicture}[font=\small, scale=0.75]
  \draw[->] (0,0) -- (6.2,0) node[right] {$t$};
  \draw[->] (0,0) -- (0,4.8) node[above] {$C_{\text{\itshape (ratio scale)}}$};
  \draw[thin, dashed] (2,0) node[below] {$t_{shock}$} -- (2,4.5);
  \draw[thick] (0.2,1.5) -- (2,1.5);
  \draw[thick] (2,1.5) .. controls (2.8,2.0) and (5.0,2.4) .. (5.9,2.4);
  \draw[thin, dash dot] (2,2.4) -- (6.2,2.4);
  \node[left] at (0,1.5) {$C^{*}$};
  \node[left] at (0,2.4) {$C^{**}$};
\end{tikzpicture}
\end{minipage}
\hfill
\begin{minipage}[t]{0.31\textwidth}
\centering
\begin{tikzpicture}[font=\small, scale=0.75]
  \draw[->] (0,0) -- (6.2,0) node[right] {$t$};
  \draw[->] (0,0) -- (0,4.8) node[above] {$Y_{\text{\itshape (ratio scale)}}$};
  \draw[thin, dashed] (2,0) node[below] {$t_{shock}$} -- (2,4.5);
  \draw[thick] (0.2,1.5) -- (2,1.5);
  \draw[thick] (2,1.5) .. controls (2.8,2.2) and (5.0,2.8) .. (5.9,2.8);
  \draw[thin, dash dot] (2,2.8) -- (6.2,2.8);
  \node[left] at (0,1.5) {$Y^{*}$};
  \node[left] at (0,2.8) {$Y^{**}$};
\end{tikzpicture}
\end{minipage}
\caption*{Figure 1}
\end{figure}

Consider the Solow Model economy studied in class:
\begin{eqnarray*}
\begin{array}{ll}
\text{Production Function} &Y_t=\overline{A} \: \left(K_t\right)^{\frac{1}{3}} \: \left(\overline{L}\right)^{\frac{2}{3}}\\[1.4ex]
\text{Capital Accumulation} &K_{t+1}=I_t+(1-\overline{d}) \: K_t \\[1.4ex]
\text{Resource Constraint} &C_t+I_t=Y_t  \\[1.4ex]
\text{Allocation of Resources} &I_t=\overline{s} \: Y_t  \\[1.4ex]
\end{array}
\end{eqnarray*}
Use this model to show the shock at time $t_{shock}$ that best describes the transitional dynamics in Figure 1. To answer this question use the Solow diagram with total output.

\subsection*{Question V: Financial Crisis (6 points)}
Suppose that the IS curve is as follows: $\widetilde{Y}_t=\frac{1}{1-b_c} \: [\overline{a}-\overline{b}_i \: (R_t-\overline{r})]-1$, $R_t=R^{ff}_t+\overline{f}$ and that the Phillips Curve is as follows: $\Delta \pi_t=\overline{v}\widetilde{Y}_t+\overline{o}_t$ . Suppose that the economy starts at its long-run equilibrium with no financial frictions ($\overline{a}=(1-b_c)$, $\overline{f}=0$, $\overline{o}_t=0$), use the IS-MP and Phillips Curve to show what happens to this economy if it experiences a negative demand shock and a concurrent increase in risk ($\overline{f}$) --- i.e. show what happens to the economy when there is a financial crisis. For this question, you may assume no monetary policy response to the crisis. 

\section{Part II: Long-Answer Questions}



\section*{Question I: The AS-AD Model}
Consider the following economy:
\begin{eqnarray*}
\begin{array}{ll}
\text{IS- Curve} & \widetilde{Y}_t=\frac{1}{1-b_c}\left[\overline{a}-\overline{b}_i \: (R_t-\overline{r})\right]-1 \\[1.4ex]
\text{Inflation} & \pi_t=\pi^e_t+\overline{v} \: \: \: \widetilde{Y}_t+\overline{o}_t, \: \: \pi^e_t=\pi_{t-1}  \\[1.4ex]
\text{Policy Rule} & R_t-\overline{r}=\overline{m}(\pi_t-\overline{\pi})  \\[1.4ex]
\end{array}
\end{eqnarray*}
Use the previous equations to answer the following questions:
\begin{enumerate}[label=(\alph*)]
\item (3 points) Derive the AD curve.
\item (6 points) With all the economic uncertainty introduced by Donald Trump's economic policies, there has been a loss in consumer and investor confidence in the US. Suppose that the economy starts at its long-run equilibrium and that the loss of confidence results in a \textbf{multi-period} negative demand shock. Using the AS-AD model, carefully explain what happens to the US economy immediately following this shock and over time.
\item (5 points) Concurrently to the negative demand shock, the increase in tariffs has resulted in a \textbf{one time} cost-pushed shock (i.e. positive inflation shock in the US, $\overline{o}_{t_{shock}}>0$).  Explain how your answer to part (b) would change if this inflation shock coincides with the first period of the demand shock. Is the initial recession milder or worse than without the inflation shock? Explain. To answer this question you must use the AS-AD graph.
\end{enumerate}


\section*{Question II: IS-MP and Phillips Curve}
Consider the following economy:
\begin{eqnarray*}
\begin{array}{ll}
\text{Expenditure:} & E_t=C_t+I_t+G_t+X_t-M_t \\[1.4ex]
\text{Consumption:} &C_t=a_c+b_c (Y_t-T_t) \\[1.4ex]
\text{Investment:} & I_t=a_i -b_i \: \: \: R_t \\[1.4ex]
\text{Imports:} & M_t=b_m \: \: \: Y_t  \\[1.4ex]
\text{Inflation:} & \pi_t=\pi_{t-1}+\overline{v} \: \: \: \widetilde{Y}_t+\overline{o}_t  \\[1.4ex]
\text{Exogenous Variables:} & G_t=\overline{G}, \: \: \: \: \: \: T_t=\overline{T}, \: \: \: \: \: \: X_t=\overline{X}
\end{array}
\end{eqnarray*}
, where $0<b_m<1$ and $0<b_c<1$. In this economy, imports depend on the level of income. Moreover, consumption depends on disposable income (income net of taxes) and not just on the level of income as we previously assumed (taxes are exogenous in this economy). Use the previous equations to answer the following questions:
\begin{enumerate}[label=(\alph*)]
\item (7 points) Solve for the equilibrium in the goods market. What is the expenditure multiplier? 
\item (7 points) Use the Keynesian Cross to carefully explain the effect of an exogenous increase in taxes ($\overline{T}$) (carefully label the graph). What is the assumption we need to make to guarantee that an equilibrium exists? Explain.
\item (6 points) Derive the IS curve using deviations from potential output as the endogenous variable.
\item (10 points) Suppose the economy is at its potential output equilibrium with a constant $2\%$ inflation rate. Assume that in the year 2000 there is a negative demand shock lasting until 2001 and decreasing the inflation rate by $1$ percentage points per year (this is a two-year recession). Assume that in 2004 the government decides to use monetary policy to increase  inflation back to its $2\%$ level by increasing output above its potential level by the same amount for two consecutive years. Using the full short-run model (IS-MP and Phillips curve), explain what happens to the economy \textbf{over time}. Be sure to include IS-MP and Phillips Curve graphs as well as graphs showing how output and inflation respond over time (carefully label all graphs).
\item (6 points) Net exports, denoted by $NX_t=\overline{X}-M_t$, model the economy's foreign trade. A positive value of $NX_t$ denotes a trade surplus and a negative value denotes a trade deficit. Suppose that trade is balanced when the economy is in its long-run equilibrium ($NX=\overline{X}-b_m \: \: \: \overline{Y}=0$), using a graph with $NX$ in the vertical axis and time on the horizontal axis show what happens to $NX_t$ from 1999 onwards.
\end{enumerate}   

\end{document}